% Conclusion (0.35 p)
\section{Conclusion}
\label{sec:conclusion}

We built a three-loop management stack for a UAV-borne 5G cell serving four
search-and-rescue robots and flew it in closed loop over real 5G software. The
most effective loop was not learned: a per-robot uplink cap driven by
round-trip time removed a silent queueing failure inside the UE, shortened C2
outages from minutes to single measurement cycles, and improved both placement
arms flown with and without it (C1). The learning layer forecasts
link-failure onset 20\,s ahead with modest skill that falls as model capacity
grows (C2), does not recover how terrain changes the effect of altitude (C3),
and, in closed loop, does not beat fixed hover or the robot centroid; made
cautious, it behaves like hover (C4). For autonomous management of such
cells, the lesson is to verify that the managed system is measured and loaded
soundly, bounding its queues and checking that its processes are alive, before
learning where to fly it. The next steps are a hardware trial of the cap and
retraining the layer on capped, watchdog-protected data.
